% Transcription draft based on Section 2 slides and video examples
\section{Differentiation}
\subsection{Maxima and minima of functions}
\begin{learningoutcomes}
\item Calculate the derivative of basic mathematical functions
\item Identify local and global maxima and minima of a function
\end{learningoutcomes}
\subsubsection{Reminder of basic differentiation}
\textbf{What is the derivative?} We have calculated the gradient of a straight line as
\[
m=\frac{\Delta y}{\Delta x}.
\]
The derivative of a function $f(x)$ is the gradient at a point (i.e. as the change in $x$ gets infinitesimally small)
\[
f'(x)=\lim_{h\to0}\frac{f(x+h)-f(x)}{h}.
\]

Another way to imagine this is demonstrated in figure \ref{fig:diff_zoom}. As we ``zoom in'' on any smooth function, it looks more like a straight line with gradient equivalent to the tangent line at that point, \textit{i.e.} the derivative.
\begin{figure}[h!]
    \centering
    \iflatexml
        \<div id="tangent-animation">\</div>
        \<script src="tangent-animation.js">\</script>
    \else
        \subfig{0.3}{Figures/Section2/plot_h_0.png}{Zoomed out, $h=1$.}
        \subfig{0.3}{Figures/Section2/plot_h_80.png}{Zooming in, $h=0.1$.}
        \subfig{0.3}{Figures/Section2/plot_h_210.png}{Zoomed in, $h=0.01$.}
    \fi
    \caption{The straight line between two points on the curve $\sin x$ gets closer to a tangent to the curve as we zoom in and reduce $h$, becoming a better approximation of the derivative at the left-hand point, which remains fixed.}
    \label{fig:diff_zoom}
\end{figure}

Some common derivatives of functions are given in table \begin{table}[ht]
    \centering
    \begin{tabular}{|c|c|} \hline
    Function $y = f(x)$ & Derivative $dy/dx = f'(x)$ \\ \hline
    $x^n$ & $nx^{n-1}$\\
    $\ln x$ & $1/x$\\
    $e^x$ & $e^x$\\
    $\sin x$ & $\cos x$\\
    $\cos x$ & $-\sin x$\\\hline
    \end{tabular}
    \caption{Some common functions and their derivatives}
    \label{tab:derivatives}
\end{table}

\textbf{The second derivative.} When a function is differentiated, the resulting derivative is also a function. It can, itself be differentiated to form the second derivative, denoted as
\[
    \der{^2y}{x^2} \qquad \mathrm{or} \qquad f''(x)
\]

\subsubsection{Local maxima and minima}

\textbf{Stationary points.} A point where $f'(x)=0$ is called a stationary point. It may represent a place where the function is a maximum or a minimum, but it could represent something else (see figure \ref{fig:stat_points}). To find stationary points, we need to solve $f'(x)=0$ for $x$, there could be
\begin{itemize}
    \item no stationary points
    \item one stationary point
    \item several stationary points.
\end{itemize}

\textbf{The first derivative test.} To use this test to classify a stationary point at $x=x_0$, 
\begin{enumerate}
    \item evaluate $f'(x)$ at two points either side of $x_0$, \emph{i.e.} $f'(x_0-h)$ and $f'(x_0+h)$ for small $h$. Make sure $h$ is small enough so that $x \pm h$ does not cross any other stationary points.
    \item Drawing the line \textbackslash for a negative value, / for a positive value and -- for the stationary point itself, sketch how the function looks around the stationary point (\textit{e.g.} a minimum will be \textbackslash -- /, and a maximum  / --\textbackslash).
\end{enumerate}

\textbf{The second derivative test.} To use this test to classify a stationary point at $x=x_0$, evaluate the second derivative $f''(x)$ at $x=x_0$.
\begin{itemize}
    \item If $f''(x_0) > 0$, the point is a minimum.
    \item If $f''(x_0) < 0$, the point is a maximum.
    \item If $f''(x_0) = 0$, the test is inconclusive. Fall back on the first derivative test. 
\end{itemize}

\begin{example}[Lecture]
Find any stationary points for the function
\[
f(x)=2x^3-9x^2+5 \, .
\]
Use the first derivative test to classify them as maxima, minima or points of inflection. Check this using the second derivative test.
\end{example}
\begin{solution}
Find the first derivative
\[
f'(x)=\blank{6x^2-18x}.
\]
Solve for $f'(x)=0$.
Stationary points are at
\[
x=\blank{0}\qquad \text{and} \qquad x=\blank{3}.
\]
First derivative test for $x=0$:\par
\vspace{4mm}
\begin{tabular}{|p{10mm}|p{10mm}|p{10mm}|p{10mm}|}\hline
    $x$ & -1 & 0 & 1\\\hline
    $f'(x)$ & \blankpar{24} & 0 & \blankpar{-12}\\\hline
    sketch & \blankpar{/} & -- & \blankpar{\textbackslash}\\\hline
\end{tabular}

\vspace{4mm}
First derivative test for $x=3$:\par
\vspace{4mm}
\begin{tabular}{|p{10mm}|p{10mm}|p{10mm}|p{10mm}|}\hline
    $x$ & 2 & 3 & 4\\\hline
    $f'(x)$ & \blankpar{-12} & 0 & \blankpar{24}\\\hline
    sketch & \blankpar{\textbackslash} & -- & \blankpar{/}\\\hline
\end{tabular}
\vspace{4mm}

Find the second derivative:
\[
f''(x)=\blank{12x-18}.
\]
\[
f''(0)=\blank{-18},\qquad f''(3)=\blank{18}.
\]
Hence $x=\blank{0}$ is a \blankpar{maximum} and $x=\blank{3}$ is a \blankpar{minimum}.
\end{solution}

\subsubsection{Maximum and minimum values over an interval}




\begin{example}[Lecture]
Find the maximum and minimum values of the function
\[
p(x)=2x^3-3x^2-36x+13
\]
on the interval $-4\le x\le 6$.
\end{example}
\begin{solution}
Find first derivative
\[
p'(x)=\blank{6x^2-6x-36}.
\]
Stationary points are at
\[
x=\blank{-2}\qquad\text{and}\qquad x=\blank{3}.
\]
Values at stationary points:
\[
f(-2)=\blank{57},\qquad f(3)=\blank{-68}.
\]
Check edges of the domain:
\[
f(-4)=\blank{29},\qquad f(6)=\blank{121}.
\]
Maximum value at $x=6$.
Minimum value at $x=3$.
\end{solution}

\begin{example}[Video]
A bioreactor produces biomass at a rate of roughly
\[
b(\omega)=3+\omega \quad [\mathrm{kg/h}],
\]
which increases with the stirring rate $\omega\geq0$ $[\mathrm{rpm}]$.
The cost of running the reactor also increases with $\omega$, roughly given by
\[
c(\omega)=72+\frac{1}{3}\omega^2 \quad [\text{\pounds}/\mathrm{h}].
\]
Find the the stirring rate that optimises production efficiency (i.e. kg of biomass per \pounds spent), and the associated production efficiency.
\end{example}
\begin{solution}
The production efficiency (in kg of biomass per \pounds spent) is given by
\[
\blank{f(\omega)=\frac{b(\omega)}{c(\omega)}=\frac{3+\omega}{72+\frac{1}{3}\omega^2}}.
\]
The optimal efficiency will be achieved when $f(\omega)$ is a maximum.
Find the stationary points $f'(\omega)=0$, using the quotient rule:
\[
\blank{f'(\omega)=\frac{\left(72+\frac{1}{3}\omega^2\right)-(3+\omega)\frac{2}{3}\omega}{\left(72+\frac{1}{3}\omega^2\right)^2}=0}.
\]
\[
\blank{72+\frac{1}{3}\omega^2-2\omega-\frac{2}{3}\omega^2=0}.
\]
Rearranging there is a stationary point when
\[
\blank{\omega^2+6\omega-216=0}.
\]
Factorising (or using the quadratic equation)
\[
\blank{(\omega+18)(\omega-12)=0}.
\]
Since $\omega\geq0$, we have an optimal stirring rate of $\omega=12$ $[\mathrm{rpm}]$.
Substituting this back into the equation for $f$, we get a production efficiency of
\[
\blank{f=\frac{3+12}{72+144/3}=\frac{15}{120}=\frac{1}{8}}\quad [\mathrm{kg}/\text{\pounds}].
\]
\end{solution}

\subsection{The Product, Quotient and Chain Rules}
\begin{learningoutcomes}
\item Identify when to use the chain, product, or quotient rules
\item Calculate the derivatives of functions using the chain, product, or quotient rules, where appropriate
\end{learningoutcomes}
\subsubsection{The chain rule}
If $f(x)=h(u)$ where $u=g(x)$ then
\[
\frac{df}{dx}=\frac{dh}{du}\frac{du}{dx}.
\]
\begin{example}[DIY]
Differentiate $(t^2+1)^3$ using the chain rule.
\end{example}
\begin{blanksol}
Let $f(t)=t^2+1$.
Then
\[
p(t)=f(t)^3.
\]
Using the chain rule,
\[
p'(t)=6t(t^2+1)^2.
\]
\end{blanksol}

\begin{example}[Video]
(a) Find the derivative of
\[
y=e^{3(1-g^2(x))}
\]
in terms of the unknown function $g(x)$ and its derivative $g'(x)$.
(b) When $g(x)=\cos x$, what is $y'$?
\end{example}
\begin{solution}
(a) First define
\[
h(x)=\blank{3(1-g^2)=3-3g^2},
\]
then
\[
h'(x)=-6\blank{gg'}.
\]
And so
\[
\blank{y'=h'e^h=-6gg'e^{3(1-g^2)}}.
\]
(b) If $g(x)=\cos x$ then $g'(x)=-\blank{\sin x}$ and therefore
\[
\blank{y'=-6gg'e^{3(1-g^2)}=6\sin x\cos x\,e^{3(1-\cos^2x)}=6\sin x\cos x\,e^{3\sin^2x}}.
\]
\end{solution}

\subsubsection{The product rule}
\begin{example}[Video]
Differentiate
\[
g(x)=x^2e^x\sin(\pi x)
\]
using the product rule.
\end{example}
\begin{solution}
Define $u(x)=x^2$, $v(x)=e^x$, and $w(x)=\sin(\pi x)$, so that $g=uvw$.
The product rule for three functions states that
\[
g'=u'vw+uv'w+uvw'.
\]
Find the derivatives of $u$, $v$, and $w$:
\[
u'=\blank{2x},\qquad v'=\blank{e^x},\qquad w'=\blank{\pi\cos(\pi x)}.
\]
And substitute them into the expression for $g'$:
\[
\blank{g'=2xe^x\sin(\pi x)+x^2e^x\sin(\pi x)+x^2e^x\pi\cos(\pi x)}.
\]
Which can be factorised as
\[
\blank{g'=xe^x\left[2\sin(\pi x)+x\sin(\pi x)+\pi x\cos(\pi x)\right]}
\]
\[
\blank{=xe^x\left[(x+2)\sin(\pi x)+\pi x\cos(\pi x)\right]}.
\]
\end{solution}

\subsubsection{The quotient rule}
\begin{example}[Video]
Differentiate $x$ with respect to $t$, where
\[
x=\frac{e^{2t}\sin(t^3)}{t}.
\]
\end{example}
\begin{solution}
Define $u(t)=e^{2t}\sin(t^3)$ and $v(t)=t$ and find their derivatives.
For $u'(t)$, we need use the product rule and chain rule:
\[
\blank{u'(t)=2e^{2t}\sin(t^3)+e^{2t}3t^2\cos(t^3)}
\]
\[
\blank{=e^{2t}\left[2\sin(t^3)+3t^2\cos(t^3)\right]}.
\]
\[
v'(t)=\blank{1}.
\]
The quotient rule states
\[
x'=\frac{u'v-v'u}{v^2}.
\]
Substituting in the expressions for $u$, $v$, $u'$, and $v'$ gives
\[
\blank{x'=\frac{e^{2t}\left[2\sin(t^3)+3t^2\cos(t^3)\right]t-e^{2t}\sin(t^3)}{t^2}}
\]
\[
\blank{=e^{2t}\left[\frac{2t-1}{t^2}\sin(t^3)+3t\cos(t^3)\right]}.
\]
\end{solution}

\subsection{Implicit, logarithmic, and parametric differentiation}
\subsubsection{Logarithmic differentiation}
\begin{example}[Video]
Find the derivative, with respect to $t$, of
\[
z=3^{t^2-1}.
\]
\end{example}
\begin{solution}
Take the natural logarithm of both sides to give
\[
\ln z=\blank{\ln\left(3^{t^2-1}\right)=(t^2-1)\ln(3)}.
\]
Differentiate both sides with respect to $t$:
\[
\blank{\frac{z'}{z}=2t\ln3}.
\]
Rearrange and substitute back in the expression for $z$ to get
\[
\blank{z'=2t\ln3\,z=2t\ln3\,3^{t^2-1}}.
\]
Alternatively, you could rearrange to write as
\[
z=\blank{e^{\ln(3)(t^2-1)}}
\]
and differentiate using the chain rule.
\end{solution}

\begin{exammistake}
Using
\[
\frac{d^2y}{dx^2}=\frac{d^2y/dt^2}{d^2x/dt^2}
\]
instead of
\[
\frac{d^2y}{dx^2}=\frac{dz/dt}{dx/dt},
\qquad z=\frac{dy}{dx}.
\]
25/26 Exam Q2(a)(i): parametric differentiation question.
\end{exammistake}
